%% kapa: exam=1 round=31 year=2020 period=1 form=A subject=civil_law subject_ko=민법 seq=36 number=36
%% source: Q-Net 감정평가사 기출문제 (cst00309) artlSeq=5207219 file=extracted/1008/5207219/A형/A형/2020년도 제31회 감정평가사 1차 시험_1교시_A형.pdf
%% conversion: kordoc 4.24.51 -> tools/gen1.py (auto); review=proofread-sonnet; flags=none
\begin{question}{36}
전세권에 관한 설명으로 옳은 것은? (다툼이 있으면 판례에 따름)
\begin{choices}
  \item 목적물의 인도는 전세권의 성립요건이다.
  \item 전세권이 존속하는 중에 전세권자는 전세권을 그대로 둔 채 전세금반환채권만을 확정적으로 양도하지 못한다.
  \item 전세목적물이 처분된 때에도 전세권을 설정한 양도인이 전세권관계에서 생기는 권리ㆍ의무의 주체이다.
  \item 전세권은 전세권설정등기의 말소등기 없이 전세기간의 만료로 당연히 소멸하지만, 전세권저당권이 설정된 때에는 그렇지 않다.
  \item 전세권저당권이 설정된 경우, 제3자의 압류 등 다른 사정이 없으면 전세권이 기간만료로 소멸한 때에 전세권설정자는 저당권자에게 전세금을 지급하여야 한다.
\end{choices}
\end{question}
\answer{2}
